\begin{abcliste}
\abc The magnetic force is the centripetal force:
\begin{align}
 e\cdot v\cdot B &= \frac{m\cdot v^2}{r} \quad\Longrightarrow\quad r = \frac{m\cdot v}{e\cdot B}
\end{align}
For $^{12}\mathrm{C}^+$ (mass $m_a = \mauO = \ma$):
\begin{align}
 r_a &= \raF \\
   &= \frac{\ma\times\v}{\nce\times\B} \\
   &= \ra \approx \result{\raP{-2}{2}}
\end{align}

\abc After half a circle an ion hits the plate at a distance $2r$ from the entrance. The distance between the two spots is
\begin{align}
 \Delta x &= 2\cdot\left(r_b-r_a\right) = \dxF \\
   &= \frac{2\times\v}{\nce\times\B}\cdot\left(\mb-\ma\right) \\
   &= \dx \approx \result{\dxP{-2}{2}}
\end{align}
(Taking $2 r_b$ or only $r_b-r_a$ instead are the typical mistakes.)

\abc Since $r = \frac{m v}{e B}$, the heavier ion moves on the larger circle: $\result{\text{$^{14}\mathrm{C}^+$}}$ hits the plate farther from the entrance ($r_b = \rbP{-2}{2}$).
\end{abcliste}
